Internet Standards enumerates design documents that
states how things work on the Internet.\todo{26/02/2026} \newline
\newline Remember the structure of the internet 
(hierarchical structure), 
given by the Backbone, then the ISPs, then organizations 
and LANs. ISPs, especially, can be divided into three 
layers, from Tier 1 to Tier 3. \newline
\newline We also have this hierarchy:
\begin{enumerate}
    \item Network Edge
    \item ?
\end{enumerate}

The number of Tier 1 ISPs is quite limited, like Verizon. 
These ones are connected to IXP (Internet Exchange Points),
which is like a switch where every company can plug its 
cables. \newline
\newline Who generates traffic? A LAN, typically. Usually, 
this LAN is able to cover only a small and localized 
area of hosts. The strucuture of a LAN is given by 
several hosts, connected to an edge router (gateway). 
This is a two-sided device, with one side being inside 
the LAN, and the other connected to the global Internet. 
A lot of protocols are used in the functioning of this 
mechanism. Remember that a single packet is encapsulated 
into different protocols. We need this encapsulation 
for several aadvantages. The upper layer protocols could 
always be the same, and we canj change lower layers 
without impacting anything. We can, for example, use a 
TCP packet over IP, and so on. \newline
\newline A local host is an host which can communicate 
directly with as local hosts. For remote hosts it's different, 
as we need to go through other routers to reach 
the other destination (so two local hosts have a direct
link between them). \newline
\newline Ethernet frames are used in this case. The
packet contains some fields like the preamble and SFD, 
which are used for synchronization. We also have destination 
and source, those contains addresses, which are called 
MAC addresses. They're 48 bits longs. We then have the 
payload and the frame clousure bits. \newline 
How do we build an Ethernet network? Specifically, 
in the real world we have several channels connected to 
a switch. So, we have many hosts connected to a single 
device, which moves packets into the Ethernet. The idea 
is that every single packet inserted into an Ethernet
network can reach any other host in the network domain, 
that's why we call this the broadcast domain. ARP, for 
example, uses broadcast messages. \newline 
Switches, however, 
do not always broadcast. Rather, it sends the packet 
on the link associated to the ethernet destination field 
addressed. Why Internet is not a large Ethernet network?
As we make a lot of use of broadcasts, which could not 
be suitable. We need to distinguish between a logical 
division of the network, and a physical one (for IP). \newline 
\textit{Part missing}. 

\section{Practical Part}

